\documentclass[11pt]{article}
\usepackage[a4paper,margin=1in]{geometry}
\usepackage[T1]{fontenc}
\usepackage{lmodern}
\usepackage{microtype}
\usepackage{amsmath,amssymb,amsthm,mathtools,mathrsfs}
\usepackage{enumitem}
\usepackage{xcolor}
\usepackage[colorlinks=true,linkcolor=blue!50!black,citecolor=blue!50!black,urlcolor=blue!50!black]{hyperref}
\usepackage[nameinlink,capitalise]{cleveref}

\newtheorem{theorem}{Theorem}[section]
\newtheorem{proposition}[theorem]{Proposition}
\newtheorem{lemma}[theorem]{Lemma}
\newtheorem{corollary}[theorem]{Corollary}
\newtheorem{question}[theorem]{Question}
\theoremstyle{definition}
\newtheorem{definition}[theorem]{Definition}
\theoremstyle{remark}
\newtheorem{remark}[theorem]{Remark}

\newcommand{\K}{\mathbb K}
\newcommand{\N}{\mathbb N}
\newcommand{\whot}{\widehat\otimes_\pi}
\newcommand{\NA}{\operatorname{NA}_\pi}
\newcommand{\FNA}{\operatorname{FNA}_\pi}
\newcommand{\INA}{\operatorname{INA}_\pi}
\newcommand{\Lin}{\mathcal L}
\newcommand{\Rea}{\operatorname{Re}}
\newcommand{\Id}{\operatorname{Id}}
\newcommand{\Cpi}{\mathscr C_\pi}
\newcommand{\cl}{\overline}

\title{Projection Descent for Norm-Attaining Projective Tensors\\
\large The gap in the published density lemma and an $L$-summand repair}
\author{Research memorandum}
\date{June 20, 2026}

\begin{document}
\maketitle

\begin{abstract}
Garc\'ia-Lirola, Guerrero-Viu, and Rueda Zoca asserted that density of projective norm-attaining tensors descends from $X\whot Y$ to $X\whot Z$ whenever $Z$ is a $1$-complemented subspace of $Y$.  The printed proof projects an ambient norm-attaining approximant and treats the projected tensor as norm-attaining.  That step is invalid: a contractive projection can destroy projective norm attainment.  Moreover, a dense-envelope construction shows that the density statement itself is false in this generality.

This memorandum gives a natural repair.  It is enough to assume that $Z$ is an $L$-summand of $Y$, equivalently that the projection $P:Y\to Z$ satisfies
\[
 \|y\|=\|Py\|+\|y-Py\|\qquad(y\in Y).
\]
In fact, for $Y=Z\oplus_1 W$, a tensor in
\[
 X\whot Y\cong (X\whot Z)\oplus_1(X\whot W)
\]
is norm-attaining if and only if both of its components are norm-attaining.  Thus the norm-attaining closure decomposes coordinatewise, and density descends exactly.

The repair is close to the natural boundary of standard projection geometry.  At the opposite $M$-summand endpoint, every tensor in $X\whot Y$ is the image of a norm-attaining tensor in $X\whot(X^*\oplus_\infty Y)$ under the canonical $M$-projection.  Hence even bicontractive $M$-projections need not preserve norm attainment.  The report distinguishes what is proved, what is merely sufficient, and what remains open.
\end{abstract}

\tableofcontents

\section{Notation and the printed statement}

All Banach spaces are over $\K\in\{\mathbb R,\mathbb C\}$.  A tensor $u\in X\whot Y$ is said to attain its projective norm if it admits a representation
\[
 u=\sum_{n=1}^\infty x_n\otimes y_n,
 \qquad
 \|u\|_\pi=\sum_{n=1}^\infty\|x_n\|\,\|y_n\|.
\]
The set of such tensors is denoted by $\NA(X\whot Y)$.  The subset of tensors admitting a finite optimal representation is denoted by $\FNA(X\whot Y)$.  We write
\[
 \Cpi(X,Y):=\cl{\NA(X\whot Y)}^{\|\cdot\|_\pi}.
\]
Finite truncation of an optimal representation gives
\begin{equation}\label{eq:closure-fna}
 \Cpi(X,Y)=\cl{\FNA(X\whot Y)}^{\|\cdot\|_\pi}.
\end{equation}

The published Lemma~4.1 of \cite{GarciaLirolaEtAl2025} states that, if
\[
 \Cpi(X,Y)=X\whot Y
\]
and $Z\subseteq Y$ is $1$-complemented, then
\[
 \Cpi(X,Z)=X\whot Z.
\]
The exact norm-attainment analogue appearing earlier as Lemma~3.1 in the same paper is correct: if every tensor in $X\whot Y$ attains its norm, then every tensor in $X\whot Z$ does so.  The passage from exact attainment to density is where the argument breaks.

\section{The precise gap in the published proof}

Let $P:Y\to Z$ be a projection with $\|P\|=1$, and put
\[
 Q=\Id_X\otimes P:X\whot Y\longrightarrow X\whot Z.
\]
The canonical inclusion $X\whot Z\hookrightarrow X\whot Y$ is isometric because $Q$ is a contractive left inverse.

\subsection{Why the earlier exact lemma is valid}

Suppose $z\in X\whot Z$ is norm-attaining when considered in $X\whot Y$.  Choose an optimal ambient representation
\[
 z=\sum_{n=1}^\infty x_n\otimes y_n,
 \qquad
 \sum_n\|x_n\|\,\|y_n\|=\|z\|_{X\whot Y}.
\]
Since $z=Qz$,
\[
 z=\sum_{n=1}^\infty x_n\otimes Py_n
\]
in $X\whot Z$.  The following chain closes at both ends:
\begin{align*}
 \|z\|_{X\whot Z}
 &\leq \sum_n\|x_n\|\,\|Py_n\|\\
 &\leq \sum_n\|x_n\|\,\|y_n\|\\
 &=\|z\|_{X\whot Y}
 =\|z\|_{X\whot Z}.
\end{align*}
Hence every inequality is an equality, and the projected representation is optimal.  This is the correct content of the published Lemma~3.1.

\subsection{The step that fails in Lemma 4.1}

For density, the paper starts with $z\in X\whot Z$ and chooses an ambient norm-attaining tensor $z'\in X\whot Y$ close to $z$.  It then says that, as in Lemma~3.1, $Qz'$ is norm-attaining.  But $z'$ need not lie in the range of $Q$, so $Qz'=z'$ is unavailable.

Indeed, from an optimal representation
\[
 z'=\sum_nx_n\otimes y_n,
 \qquad
 \|z'\|_\pi=\sum_n\|x_n\|\,\|y_n\|,
\]
one obtains only
\begin{equation}\label{eq:broken-chain}
 \|Qz'\|_\pi
 \leq\sum_n\|x_n\|\,\|Py_n\|
 \leq\sum_n\|x_n\|\,\|y_n\|
 =\|z'\|_\pi.
\end{equation}
There is no equality at the right endpoint involving $\|Qz'\|_\pi$, and the first inequality in \eqref{eq:broken-chain} may be strict.  The proof silently assumes precisely the equality that has to be established.

The defect can be written as
\[
 \delta_P(z';(x_n,y_n))
 :=\sum_n\|x_n\|\,\|Py_n\|-\|Qz'\|_\pi\geq0.
\]
Lemma~3.1 forces this defect to vanish because the tensor itself belongs to the projected range.  For a general approximant there is no such mechanism.

\subsection{Projection can genuinely destroy norm attainment}

The missing assertion is not merely unproved; it is false in a particularly strong form.

\begin{theorem}[Universal $M$-projection lift]\label{thm:M-lift}
Let $X,Y$ be Banach spaces and let $z\in X\whot Y$.  For every representation
\[
 z=\sum_{n=1}^\infty x_n\otimes y_n,
 \qquad
 \sum_n\|x_n\|\,\|y_n\|<\infty,
\]
there exists
\[
 \widetilde z\in\NA\bigl(X\whot(X^*\oplus_\infty Y)\bigr)
\]
such that
\[
 (\Id_X\otimes R)\widetilde z=z,
\]
where $R:X^*\oplus_\infty Y\to Y$ is the second-coordinate projection.  If the chosen representation of $z$ is finite, then $\widetilde z$ belongs to $\FNA$.
\end{theorem}

\begin{proof}
Omit zero terms.  For each $n$, choose $x_n^*\in S_{X^*}$ with
\[
 x_n^*(x_n)=\|x_n\|.
\]
Set
\[
 f_n=\|y_n\|x_n^*,
 \qquad
 \widetilde z=\sum_nx_n\otimes(f_n,y_n).
\]
Since
\[
 \|(f_n,y_n)\|_{X^*\oplus_\infty Y}=\|y_n\|,
\]
the series converges projectively.  Define the trace form
\[
 \tau(x,(f,y))=f(x).
\]
It has norm one, and
\[
 \tau\bigl(x_n,(f_n,y_n)\bigr)
 =\|x_n\|\,\|y_n\|.
\]
Therefore
\[
 \|\widetilde z\|_\pi
 \geq \tau(\widetilde z)
 =\sum_n\|x_n\|\,\|y_n\|
 \geq\|\widetilde z\|_\pi,
\]
so the displayed representation is optimal.  The projection identity is immediate.  Finally,
\[
 \|(f,y)\|_\infty
 =\max\{\|(I-R)(f,y)\|,\|R(f,y)\|\},
\]
so $R$ is an $M$-projection, in particular a bicontractive projection.
\end{proof}

\begin{corollary}\label{cor:M-destroys}
There are a bicontractive $M$-projection $R:F\to Y$ and a norm-attaining tensor $v\in X\whot F$ such that $(\Id_X\otimes R)v$ is not norm-attaining.  One may take $v$ to have a finite optimal representation while its image is a finite-rank non-norm-attaining tensor.
\end{corollary}

\begin{proof}
Non-norm-attaining tensors exist, and finite-rank examples are given in \cite[Proposition 5.5]{DantasEtAl2022}.  Apply \cref{thm:M-lift} to such a tensor and to a finite representation in the finite-rank case.
\end{proof}

Thus neither contractivity, bicontractivity, nor the stronger $M$-projection identity repairs the specific projection step in the printed proof.

\section{Why the density statement itself is false}

The preceding theorem disproves preservation of individual norm-attaining tensors, but by itself it does not disprove density descent.  We record a self-contained dense-envelope construction that does.

\subsection{A trace-face correction lemma}

Let $E,H$ be Banach spaces and set $G=E^*\oplus_\infty H$.  Define
\[
 \tau_E(e,(e^*,h))=e^*(e).
\]

\begin{lemma}[Trace-face correction]\label{lem:trace-face}
If $u\in E\whot G$ satisfies
\[
 \Rea\tau_E(u)=\|u\|_\pi,
\]
then $u\in\Cpi(E,G)$.  More precisely, $u$ is approximable by tensors with finite optimal representations supported by $\tau_E$.
\end{lemma}

\begin{proof}
Normalize $\|u\|_\pi=1$.  Fix $k\in\N$ and let $\eta_k>0$ be a scalar Bishop--Phelps--Bollob\'as threshold for error $1/k$.  Choose a projective representation
\[
 u=\sum_n\lambda_{k,n}e_{k,n}\otimes(f_{k,n},h_{k,n})
\]
with
\[
 e_{k,n}\in S_E,
 \quad
 (f_{k,n},h_{k,n})\in S_G,
 \quad
 \sum_n\lambda_{k,n}<1+\frac{\eta_k}{k}.
\]
The trace equality gives
\[
 \sum_n\lambda_{k,n}
 \bigl(1-\Rea f_{k,n}(e_{k,n})\bigr)
 <\frac{\eta_k}{k}.
\]
Let
\[
 A_k=\{n:\Rea f_{k,n}(e_{k,n})>1-\eta_k\}.
\]
Then $\sum_{n\notin A_k}\lambda_{k,n}<1/k$.  For $n\in A_k$, the Bishop--Phelps--Bollob\'as theorem gives $\widetilde e_{k,n}\in S_E$ and $\widetilde f_{k,n}\in S_{E^*}$ such that
\[
 \widetilde f_{k,n}(\widetilde e_{k,n})=1,
 \quad
 \|e_{k,n}-\widetilde e_{k,n}\|<\frac1k,
 \quad
 \|f_{k,n}-\widetilde f_{k,n}\|<\frac1k.
\]
Since $\|h_{k,n}\|\leq1$,
\[
 \|(\widetilde f_{k,n},h_{k,n})\|=1.
\]
The tensor
\[
 v_k=\sum_{n\in A_k}\lambda_{k,n}
 \widetilde e_{k,n}\otimes(\widetilde f_{k,n},h_{k,n})
\]
is norm-attaining, supported by $\tau_E$, and $\|u-v_k\|_\pi\to0$.  Finite partial sums of $v_k$ are finitely norm-attaining, and a diagonal selection proves the result.
\end{proof}

\subsection{Universal graph stabilization}

For fixed $X,Y$, let
\[
 \Gamma_Y=B_{\Lin(X,Y^*)},
 \qquad
 \mathcal U_X(Y)=Y\oplus_\infty\ell_\infty(\Gamma_Y,X^*).
\]
For $B\in\Gamma_Y$, write
\[
 (B^{\mathsf t}y)(x)=B(x)(y).
\]
Define
\[
 j_Y(y)=\bigl(y,(B^{\mathsf t}y)_{B\in\Gamma_Y}\bigr).
\]
The map $j_Y$ is an isometry, and its range is $1$-complemented by
\[
 Q_Y\bigl(y,(f_B)_B\bigr)=j_Y(y).
\]

\begin{proposition}[Universal graph stabilization]\label{prop:graph}
For every $u\in X\whot Y$,
\[
 (\Id_X\otimes j_Y)u\in\Cpi\bigl(X,\mathcal U_X(Y)\bigr).
\]
\end{proposition}

\begin{proof}
Choose $B_0\in S_{\Lin(X,Y^*)}$ with
\[
 \Rea B_0(u)=\|u\|_\pi.
\]
The $B_0$ coordinate of $\mathcal U_X(Y)$ is a copy of $X^*$.  Its coordinate trace has norm one and evaluates $(\Id_X\otimes j_Y)u$ as $B_0(u)$.  The graph embedding is isometric, so the trace value equals the projective norm.  Apply \cref{lem:trace-face}.
\end{proof}

\begin{theorem}[Dense envelope]\label{thm:dense-envelope}
For every pair $X,Y$, there is a Banach space $Y_\infty$ containing $Y$ isometrically as a $1$-complemented subspace such that
\[
 \Cpi(X,Y_\infty)=X\whot Y_\infty.
\]
\end{theorem}

\begin{proof}
Set $Y_0=Y$ and define inductively
\[
 Y_{n+1}=\mathcal U_X(Y_n),
\]
identifying $Y_n$ with its graph image in $Y_{n+1}$.  Every inclusion is isometric and has a contractive retraction.  Let $Y_\infty$ be the completion of the increasing union.  Compatibility of the retractions makes each $Y_n$, especially $Y_0$, $1$-complemented in $Y_\infty$.

By \cref{prop:graph}, every tensor in $X\whot Y_n$ belongs to $\Cpi(X,Y_{n+1})$.  Since $Y_{n+1}$ is $1$-complemented in $Y_\infty$, finite optimal representations in $X\whot Y_{n+1}$ remain optimal in the ambient tensor product.  Hence
\[
 X\whot Y_n\subseteq\Cpi(X,Y_\infty)
 \qquad(n\in\N).
\]
The union of the left-hand spaces is dense in $X\whot Y_\infty$, so the closed set $\Cpi(X,Y_\infty)$ is the whole space.
\end{proof}

\begin{corollary}[Failure of the printed Lemma 4.1]\label{cor:printed-false}
There are Banach spaces $X,Y,Z$ such that $Y$ is an isometric $1$-complemented subspace of $Z$,
\[
 \Cpi(X,Z)=X\whot Z,
\]
but
\[
 \Cpi(X,Y)\neq X\whot Y.
\]
Consequently, Lemma~4.1 of \cite{GarciaLirolaEtAl2025} is false as printed.
\end{corollary}

\begin{proof}
Choose a pair $X,Y$ with non-dense norm-attaining tensors, whose existence is proved in \cite[Theorem 5.1]{DantasEtAl2022}, and apply \cref{thm:dense-envelope}.  Take $Z=Y_\infty$.
\end{proof}

The construction is new and has not undergone independent peer review.  The gap in the printed proof, however, is independent of the construction and is already witnessed by \cref{thm:M-lift}.

\section{The $L$-summand repair}

\begin{definition}
A projection $P:Y\to Y$ is an \emph{$L$-projection} if
\[
 \|y\|=\|Py\|+\|y-Py\|
 \qquad(y\in Y).
\]
Its range is called an \emph{$L$-summand}.  Equivalently, if $Z=P(Y)$ and $W=\ker P$, then
\[
 Y=Z\oplus_1 W
\]
isometrically.  An $M$-projection is defined by replacing the sum with the maximum.
\end{definition}

The $L$-identity supplies exactly the equality missing from \eqref{eq:broken-chain}.

\begin{lemma}[Projective distributivity over an $\ell_1$-sum]\label{lem:l1-distribution}
For Banach spaces $X,Z,W$, the canonical map
\[
 \Phi:X\whot(Z\oplus_1W)
 \longrightarrow
 (X\whot Z)\oplus_1(X\whot W)
\]
defined on elementary tensors by
\[
 \Phi\bigl(x\otimes(z,w)\bigr)=(x\otimes z,x\otimes w)
\]
is an onto linear isometry.
\end{lemma}

\begin{proof}
For an elementary tensor,
\[
 \|x\otimes z\|_\pi+\|x\otimes w\|_\pi
 =\|x\|(\|z\|+\|w\|)
 =\|x\otimes(z,w)\|_\pi.
\]
The projective definition shows that $\Phi$ is contractive.  Conversely, the map
\[
 (u,v)\longmapsto (\Id_X\otimes i_Z)u+(\Id_X\otimes i_W)v
\]
is contractive from $(X\whot Z)\oplus_1(X\whot W)$ into $X\whot(Z\oplus_1W)$ and is inverse to $\Phi$ on a dense algebraic subspace.  Hence both maps extend as inverse isometries.
\end{proof}

\begin{theorem}[Exact decomposition of norm attainment]\label{thm:l1-NA}
Under the isometry in \cref{lem:l1-distribution}, let $u$ correspond to $(u_Z,u_W)$.  Then
\[
 u\in\NA\bigl(X\whot(Z\oplus_1W)\bigr)
\quad\Longleftrightarrow\quad
 u_Z\in\NA(X\whot Z)
 \text{ and }
 u_W\in\NA(X\whot W).
\]
The same equivalence holds with $\FNA$ in place of $\NA$.
\end{theorem}

\begin{proof}
Assume first that
\[
 u=\sum_{n=1}^\infty x_n\otimes(z_n,w_n)
\]
is an optimal representation.  Put
\[
 A=\sum_n\|x_n\|\,\|z_n\|,
 \qquad
 C=\sum_n\|x_n\|\,\|w_n\|.
\]
Because the norm on $Z\oplus_1W$ is additive,
\[
 A+C
 =\sum_n\|x_n\|(\|z_n\|+\|w_n\|)
 =\|u\|_\pi.
\]
By \cref{lem:l1-distribution},
\[
 \|u\|_\pi=\|u_Z\|_\pi+\|u_W\|_\pi.
\]
The projected series give
\[
 \|u_Z\|_\pi\leq A,
 \qquad
 \|u_W\|_\pi\leq C.
\]
The sums of the two sides are equal, so both inequalities are equalities.  Thus both projected representations are optimal.  Finiteness is preserved.

Conversely, concatenate optimal representations of $u_Z$ and $u_W$, embedding the first in the $Z$ coordinate and the second in the $W$ coordinate.  The total coefficient sum is
\[
 \|u_Z\|_\pi+\|u_W\|_\pi=\|u\|_\pi,
\]
so the resulting representation of $u$ is optimal.  Again, finite representations remain finite.
\end{proof}

\begin{corollary}[Closure decomposition]\label{cor:closure-decomp}
Under the same identification,
\[
 \Cpi(X,Z\oplus_1W)
 =\Cpi(X,Z)\oplus_1\Cpi(X,W).
\]
In particular,
\[
 \Cpi(X,Z\oplus_1W)=X\whot(Z\oplus_1W)
\]
if and only if
\[
 \Cpi(X,Z)=X\whot Z
 \quad\text{and}\quad
 \Cpi(X,W)=X\whot W.
\]
\end{corollary}

\begin{proof}
By \cref{thm:l1-NA}, the norm-attaining set corresponds exactly to the Cartesian product of the two coordinate norm-attaining sets.  Taking closure in an $\ell_1$-sum gives the product of the two closures.
\end{proof}

\begin{theorem}[Corrected density-descent lemma]\label{thm:corrected-lemma}
Let $X,Y$ be Banach spaces and let $Z$ be an $L$-summand of $Y$.  If
\[
 \Cpi(X,Y)=X\whot Y,
\]
then
\[
 \Cpi(X,Z)=X\whot Z.
\]
More strongly, if $P:Y\to Z$ is the $L$-projection, then
\[
 (\Id_X\otimes P)
 \NA(X\whot Y)
 =\NA(X\whot Z),
\]
and the same equality holds for $\FNA$.
\end{theorem}

\begin{proof}
Write $Y=Z\oplus_1\ker P$ and apply \cref{thm:l1-NA,cor:closure-decomp}.  The reverse inclusion in the displayed equality follows because every intrinsic optimal representation remains optimal under the canonical isometric inclusion into the ambient space.
\end{proof}

This is the cleanest direct repair of the printed argument: replace ``$1$-complemented'' by ``an $L$-summand.''

\subsection{Arbitrary $\ell_1$-sums}

The same proof gives a useful extension.  If $(Y_i)_{i\in I}$ is any family, then
\[
 X\whot\left(\bigoplus_{i\in I}Y_i\right)_1
 \cong
 \left(\bigoplus_{i\in I}X\whot Y_i\right)_1,
\]
and a tensor $(u_i)_i$ is norm-attaining if and only if every nonzero component $u_i$ is norm-attaining.  Consequently,
\[
 \Cpi\left(X,\left(\bigoplus_{i\in I}Y_i\right)_1\right)
 =
 \left(\bigoplus_{i\in I}\Cpi(X,Y_i)\right)_1.
\]
Every element of an $\ell_1(I)$-sum has countable support, so concatenating the coordinate optimal representations causes no set-theoretic difficulty.

\subsection{Bidual and integral versions}

A Banach space $Y$ is $L$-embedded if its canonical copy is an $L$-summand in $Y^{**}$.  Therefore \cref{thm:corrected-lemma} gives the valid bidual replacement
\[
 \Cpi(X,Y^{**})=X\whot Y^{**}
 \quad\Longrightarrow\quad
 \Cpi(X,Y)=X\whot Y
\]
for $L$-embedded $Y$.

The 2026 preprint \cite{AliagaEtAl2026} proves that density of classical norm-attaining tensors is equivalent to density of integral norm-attaining tensors.  Combining that equivalence with \cref{thm:corrected-lemma} yields the corresponding corrected version of its Corollary~2.14: density of integral norm-attaining tensors descends through $L$-summands.  The statement with an arbitrary $1$-complemented subspace is contradicted by \cref{cor:printed-false}.

\section{How restrictive is the repair?}

The phrase ``least restrictive'' has to be interpreted carefully.  There cannot be a unique weakest geometric hypothesis on $P$: pair-specific assumptions on $X$, $Y$, and $Z$ may imply intrinsic density for reasons unrelated to the projection.  There is, however, a precise weakest condition for the published proof strategy itself:
\[
 (\Id_X\otimes P)\NA(X\whot Y)
 \subseteq\NA(X\whot Z).
\]
This condition is tautological rather than structural.  The $L$-summand assumption is a standard, checkable geometric condition that forces it, and it does so with an exact coordinate characterization.

The comparison with the two classical absolute-sum endpoints is sharp for the projection mechanism:
\begin{itemize}[leftmargin=2em]
\item For an $\ell_1$-decomposition, every norm-attaining tensor decomposes into norm-attaining coordinates by \cref{thm:l1-NA}.
\item For an $\ell_\infty$-decomposition, the phenomenon reverses maximally: by \cref{thm:M-lift}, the coordinate $M$-projection maps ambient norm-attaining tensors onto every tensor of the target, including non-norm-attaining ones.
\end{itemize}
Thus ``bicontractively complemented'' or ``an $M$-summand'' is not enough to validate the sentence used in the published proof.  This does not by itself decide whether the density conclusion might hold for all $M$-summands by some different argument; that narrower question appears to remain open.

For a nonzero factor $X$, the tensor projection $\Id_X\otimes P$ is an $L$-projection if and only if $P$ is an $L$-projection.  The nontrivial implication follows by testing the $L$-identity on $x\otimes y$ with $x\in S_X$:
\[
 \|y\|
 =\|x\otimes y\|_\pi
 =\|x\otimes Py\|_\pi+\|x\otimes(y-Py)\|_\pi
 =\|Py\|+\|y-Py\|.
\]
So the repair cannot be weakened merely by moving the same $L$-identity from the factor to the tensor product.

\section{Conclusions}

The status can be summarized precisely.

\begin{enumerate}[leftmargin=2em]
\item The exact all-tensors descent Lemma~3.1 of \cite{GarciaLirolaEtAl2025} is correct for arbitrary $1$-complemented subspaces.
\item The density proof for Lemma~4.1 has a definite gap at the claim that the projection of an ambient norm-attaining approximant is norm-attaining.
\item That preservation claim is false even for bicontractive $M$-projections, by \cref{thm:M-lift}.
\item The density statement itself is false for arbitrary $1$-complemented subspaces, by the dense-envelope construction.
\item A natural corrected theorem is obtained by requiring the subspace to be an $L$-summand.  Under this assumption, norm attainment and its closure decompose exactly coordinatewise.
\end{enumerate}

The $L$-summand theorem is therefore a rigorous and near-sharp repair in the standard geometry of contractive projections.  It should not be advertised as the logically weakest possible hypothesis, but it is the next natural condition after $1$-complementability that restores the missing equality and makes the original proof strategy valid.

\begin{thebibliography}{9}

\bibitem{AliagaEtAl2026}
R.~J. Aliaga, S.~Dantas, J.~Guerrero-Viu, M.~Jung, and O.~Rold\'an,
\emph{Integral representations of projective norm-attaining tensors},
arXiv:2602.23796, 2026.

\bibitem{Bollobas1970}
B.~Bollob\'as,
\emph{An extension to the theorem of Bishop and Phelps},
Bull. London Math. Soc. 2 (1970), 181--182.

\bibitem{DantasEtAl2022}
S.~Dantas, M.~Jung, O.~Rold\'an, and A.~Rueda Zoca,
\emph{Norm-attaining tensors and nuclear operators},
Mediterr. J. Math. 19 (2022), article 38; arXiv:2006.09871.

\bibitem{GarciaLirolaEtAl2025}
L.~C. Garc\'ia-Lirola, J.~Guerrero-Viu, and A.~Rueda Zoca,
\emph{Projective tensor products where every element is norm-attaining},
Banach J. Math. Anal. 19 (2025), article 19,
\href{https://doi.org/10.1007/s43037-024-00400-7}{doi:10.1007/s43037-024-00400-7}.

\bibitem{HWW}
P.~Harmand, D.~Werner, and W.~Werner,
\emph{$M$-Ideals in Banach Spaces and Banach Algebras},
Lecture Notes in Mathematics 1547, Springer, 1993.

\bibitem{Ryan}
R.~A. Ryan,
\emph{Introduction to Tensor Products of Banach Spaces},
Springer Monographs in Mathematics, Springer, 2002.

\end{thebibliography}

\end{document}
